\chapter{Worked examples: discrete-logarithm schemes and protocols}

These modules work over an abstract cyclic group presented with explicit
exponents, in contrast to the pairing-group formulation used in the ElGamal and
Diffie--Hellman chapters. Public-key encryption
follows Rosulek, \emph{The Joy of Cryptography}, Chapter~15; oblivious transfer,
the Chaum--Pedersen protocol and coin tossing go beyond the book.

\section{Cyclic groups and DDH}

\begin{definition}[Cyclic group with exponents; DDH games]
  \label{def:cyclicgroup-ddh}
  \uses{def:spcomp, def:advantage}
  \lean{CatCryptCore.Examples.CyclicGroupDDH.CyclicGroup, CatCryptCore.Examples.CyclicGroupDDH.DDH_real, CatCryptCore.Examples.CyclicGroupDDH.DDH_Advantage}
  \leanok
  A finite group with a bijective exponentiation $\mathsf{pow} : \mathsf{Exp} \to G$
  of the generator; the real game outputs $(g^a, g^b, g^{ab})$, the ideal game
  $(g^a, g^b, g^c)$, and the DDH advantage is the distance between them.
\end{definition}

\begin{definition}[Prime-order group parameters]
  \label{def:groupparam}
  \lean{CatCrypt.Examples.GroupParam}
  \leanok
  A cyclic group of prime order with its scalar field, packaged as one structure
  whose fields record the group and field axioms.
\end{definition}

\section{ElGamal under DDH}

\begin{theorem}[Correctness and IND-CPA security]
  \label{thm:elgamal-cyclic}
  \uses{def:cyclicgroup-ddh, thm:advantageA-triangle}
  \lean{CatCryptCore.Examples.ElGamalDDH.ElGamalScheme, CatCryptCore.Examples.ElGamalDDH.elgamal_correct, CatCryptCore.Examples.ElGamalDDH.Reduction, CatCryptCore.Examples.ElGamalDDH.elgamal_indcpa_security}
  \leanok
  Decryption inverts encryption, and for messages $m_0, m_1$ and every adversary
  $A$ the IND-CPA advantage is at most
  $\Adv^{\mathrm{DDH}}(R_{m_0}(A)) + \Adv^{\mathrm{DDH}}(R_{m_1}(A))$ for the
  explicit reduction $R_m$, which embeds a DDH triple as public key and
  ciphertext.
\end{theorem}

\begin{theorem}[Concrete UC emulation]
  \label{thm:elgamal-uc}
  \uses{thm:elgamal-cyclic}
  \lean{CatCryptCore.Examples.ElGamalUCConcrete.elgamalGameAdvBoundC, CatCryptCore.Examples.ElGamalUCConcrete.elgamal_uc_concrete, CatCryptCore.Examples.ElGamalUCConcrete.elgamal_uc_concrete_apply}
  \leanok
  The ElGamal game pair, indexed by the message pair, is a concrete UC emulation:
  for every adversary $A$ and environment $Z$ the gap is at most the two DDH
  advantages of the reductions applied to the composite of $A$ and $Z$.
\end{theorem}

\section{Hashed ElGamal}

\begin{theorem}[Correctness and the DDH reduction]
  \label{thm:hashed-elgamal}
  \uses{def:cyclicgroup-ddh, thm:advantageA-triangle}
  \lean{CatCryptCore.Examples.HashedElGamal.heg_correct_bool, CatCryptCore.Examples.HashedElGamal.heg_ddh_reduction, CatCryptCore.Examples.HashedElGamal.heg_indcpa_le_ddh, CatCryptCore.Examples.HashedElGamal.heg_perfect_hash_security}
  \leanok
  Encryption masks a bit with $H(\mathit{pk}^r)$ and decryption removes the mask.
  The distance between the real games for $m_0$ and $m_1$ is at most two DDH
  advantages of explicit reductions, under the hypothesis that the reductions run
  on the ideal DDH triple are independent of the message. Separately, when the
  mask is a uniform bit the two games have advantage exactly $0$.
\end{theorem}

\section{Oblivious transfer}

The Naor--Pinkas 1-out-of-2 oblivious transfer.

\begin{theorem}[Sender privacy for one encryption]
  \label{thm:ot-sender}
  \uses{def:cyclicgroup-ddh}
  \lean{CatCrypt.Examples.OT.CyclicGroupRing, CatCrypt.Examples.OT.otEnc_message_independent}
  \leanok
  When the exponent $d$ at the non-DDH position satisfies $d \neq ab$, the map
  $(s, r) \mapsto (as + r,\ ds + br)$ is a bijection, so the sender's encryption
  at that position has the same distribution for every message.
\end{theorem}

\begin{theorem}[Sender privacy for the transcript]
  \label{thm:ot-sender-game}
  \uses{thm:ot-sender}
  \lean{CatCrypt.Examples.OT.advantageA_sample_bind_le_of_eq_off, CatCrypt.Examples.OT.ot_sender_secure}
  \leanok
  For every choice bit and every distinguisher, the real transcript and the ideal
  transcript, in which the unchosen message is replaced by the identity, are at
  advantage at most $1/|\mathsf{Exp}|$. The two games agree on every receiver
  sample with $c \neq ab$; the bound is the probability of the event $c = ab$.
\end{theorem}

\begin{theorem}[Receiver privacy]
  \label{thm:ot-receiver}
  \uses{def:cyclicgroup-ddh, thm:advantageA-triangle}
  \lean{CatCrypt.Examples.OT.receiverView, CatCrypt.Examples.OT.receiver_security}
  \leanok
  The advantage in distinguishing the receiver's message for choice bit $0$ from
  that for choice bit $1$ is at most the sum of two DDH advantages of explicit
  reductions.
\end{theorem}

\section{The Chaum--Pedersen protocol}

A $\Sigma$-protocol for equality of discrete logarithms, $u = g^x$ and $v = h^x$,
built on the $\Sigma$-protocol interface of the $\Sigma$-protocol chapter.

\begin{theorem}[Completeness, SHVZK and special soundness]
  \label{thm:chaum-pedersen}
  \uses{def:groupparam, thm:sigma-complete}
  \lean{CatCryptCore.Examples.ChaumPedersen.chaumPedersenSigma, CatCryptCore.Examples.ChaumPedersen.chaumPedersen_complete, CatCryptCore.Examples.ChaumPedersen.chaumPedersen_SHVZK_visible, CatCryptCore.Examples.ChaumPedersen.chaumPedersen_special_soundness}
  \leanok
  Honest transcripts verify; for each challenge the real and simulated
  transcripts are equal in distribution; two accepting transcripts with one
  commitment and distinct challenges yield the common exponent.
\end{theorem}

\begin{theorem}[UC security]
  \label{thm:chaum-pedersen-uc}
  \uses{thm:chaum-pedersen}
  \lean{CatCryptCore.Examples.ChaumPedersen.chaumPedersen_uc_secure}
  \leanok
  For a statement in the relation, the real protocol UC-emulates the ideal one
  with error $0$.
\end{theorem}

\section{Coin tossing}

\begin{theorem}[Two-party coin toss is uniform]
  \label{thm:cointoss}
  \uses{def:sdistr-uniform}
  \lean{CatCryptCore.Examples.CoinToss.coinTossReal, CatCryptCore.Examples.CoinToss.coinTossIdeal, CatCryptCore.Examples.CoinToss.add_uniform_eq_uniform, CatCryptCore.Examples.CoinToss.coinToss_eq}
  \leanok
  The sum of two independent uniform elements of $\ZZ/p$ is uniform: the real
  protocol and the ideal functionality are equal as computations.
\end{theorem}
